% Section 6: Contextual Benchmark Workloads
% Describes the C1--C10 workload taxonomy and fairness methodology.

\section{Contextual Benchmark Workloads}
\label{sec:contextual-workloads}

Standard RDF benchmark suites (e.g.,~SP$^2$Bench, LUBM, WatDiv) focus on
relational-style queries -- selective lookups, multi-way joins, aggregation --
that map naturally to the relational algebra underlying SPARQL engines.
SheafDB's advantage lies precisely where the relational model struggles:
contexts that span multiple facts, high-arity relations that resist
reification, temporal intervals that require range-indexed access,
provenance DAGs, and global consistency constraints.

To evaluate these dimensions systematically, we introduce ten contextual
benchmark workloads (C1--C10), organized into six categories:

\begin{itemize}
    \item \textbf{Contextual Neighborhood (C1--C2):} Retrieve facts sharing
    entities with a seed fact within a bounded radius; enumerate paths
    connecting two entities across the fact graph.
    \item \textbf{High Arity (C3--C4):} Look up and join facts whose
    relations carry 3--30 arguments, exercising SheafDB's native tuple
    store against KG reification overhead.
    \item \textbf{Temporal (C5--C6):} Interval overlap queries and
    sliding-window aggregation over facts with valid-time ranges.
    \item \textbf{Provenance (C7):} Trace derivation chains through
    source--transformation--confidence DAGs.
    \item \textbf{Consistency \& Global Sections (C8--C9):} Verify
    cocycle conditions on overlapping context assignments; construct
    global sections from compatible local data.
    \item \textbf{Mixed Flagship (C10):} A multi-stage pipeline combining
    neighborhood, temporal, provenance, and consistency checks --- the
    primary evaluation workload for the SheafDB paper.
\end{itemize}

\subsection{Fairness Contract}
\label{sec:fairness}

Every workload defines a \emph{semantic question} that both engines answer.
SheafDB-native operations (neighborhood, consistency, global-section)
document their KG-equivalent simulation via SPARQL.  Three fairness gates
are enforced:

\begin{enumerate}
    \item \textbf{Executability:} Both engines can produce a result (possibly
    empty) for every workload.
    \item \textbf{Equivalence:} Result sets are compared as sets of frozen
    dictionaries; row-wise equality up to column ordering.
    \item \textbf{Documentation:} Any operator available in one engine but
    not the other is recorded with a simulation note in the fairness report.
\end{enumerate}

Section~\ref{sec:results} reports the per-workload speedup and equivalence
status.
